\chapter{تبلیغات برخط در وب \texorpdfstring{\enparen{Online Advertising \& AdWords}}{(Online Advertising \& AdWords)}}
\label{ch:advertising}

\section{مقدمه: الگوریتم‌های برخط و اقتصاد موتورهای جستجو}
درآمد اصلی غول‌های فناوری وب (به‌ویژه گوگل و متا) از تبلیغات هدفمند اینترنتی حاصل می‌شود. تفاوت بنیادین تبلیغات سنتی (نظیر بیلبورد یا آگهی روزنامه) با تبلیغات وب در دو عامل خلاصه می‌شود:
\begin{enumerate}
    \item \textbf{هدفمندی مستقیم}\footnote{\lr{Targeted Advertising}}: نمایش تبلیغ متناسب با نیاز دقیق لحظه‌ای کاربر بر اساس پرس‌وجوی جستجو\footnote{\lr{Search Query}}.
    \item \textbf{الگوریتم‌های برخط}\footnote{\lr{Online Algorithms}}: ورودی‌های مسئله (پرس‌وجوهای کاربران) یکی‌یکی در طول زمان وارد می‌شوند و الگوریتم باید بلافاصله، بدون اطلاع از آینده و به صورت بازگشت‌ناپذیر تصمیم بگیرد که کدام آگهی را نمایش دهد.
\end{enumerate}

در الگوریتم‌های سنتی، کل ورودی از ابتدا معلوم است (الگوریتم آفلاین)، اما در الگوریتم‌های برخط، تصمیم‌گیری در شرایط نااطمینانی صورت می‌پذیرد.

\begin{definition}[نسبت رقابتی \enparen{Competitive Ratio}]
کارایی یک الگوریتم برخط در مقایسه با یک الگوریتم بهینه آفلاین (که با برخورداری از ویژگی «پیش‌گویی»، تمام آینده را می‌داند) با معیاری به نام \textbf{نسبت رقابتی}\footnote{\lr{Competitive Ratio}} ارزیابی می‌شود:
\begin{equation}
\text{CR} = \min_{\text{ورودی‌های ممکن } I} \frac{\text{سود الگوریتم برخط }(I)}{\text{سود الگوریتم بهینه آفلاین }(I)}
\end{equation}
نسبت رقابتی مقداری بین ۰ و ۱ است. هرچه این نسبت به ۱ نزدیک‌تر باشد، عملکرد الگوریتم برخط بهتر است.
\end{definition}

\section{مسئله تطابق دوبخشی برخط \texorpdfstring{\enparen{Online Bipartite Matching}}{(Online Bipartite Matching)}}
مسئله نمایش تبلیغات را می‌توان در قالب یک گراف دوبخشی $G = (U, V, E)$ مدل‌سازی کرد:
\begin{itemize}
    \item بخش $U$: تبلیغ‌دهندگان\footnote{\lr{Advertisers}} که از قبل در سیستم ثبت‌نام کرده‌اند و ثابت هستند.
    \item بخش $V$: پرس‌وجوهای کاربران\footnote{\lr{Queries}} که یکی‌یکی به صورت آنلاین وارد می‌شوند.
    \item یال $(u, v)$: نشان می‌دهد که تبلیغ‌دهنده $u$ تمایل به نمایش آگهی برای پرس‌وجوی $v$ دارد.
\end{itemize}

هدف: یافتن بزرگ‌ترین تطابق\footnote{\lr{Maximum Matching}}؛ یعنی بیشترین تعداد یال‌های بدون رأس مشترک.

\subsection{الگوریتم حریصانه \texorpdfstring{\enparen{Greedy Algorithm}}{(Greedy Algorithm)} و تحلیل نسبت رقابتی}
ساده‌ترین استراتژی برخط، الگوریتم حریصانه است: به محض ورود یک پرس‌وجو $v$، آن را به هر تبلیغ‌دهنده متصل به آن که تاکنون آگهی به او تخصیص داده نشده، اختصاص می‌دهیم.

\begin{theorem}[کران نسبت رقابتی الگوریتم حریصانه]
نسبت رقابتی الگوریتم حریصانه برای تطابق دوبخشی برخط حداقل برابر با $\frac{1}{2}$ $(50\%)$ است و این کران تنگ است.
\end{theorem}

\textbf{اثبات قضیه:}
فرض کنید تطابق تولیدشده توسط الگوریتم حریصانه $M_{\text{greedy}}$ و تطابق بهینه آفلاین $M_{\text{opt}}$ باشد.
می‌دانیم هر یال $(u, v) \in M_{\text{opt}}$ باید حداقل با یکی از یال‌های $M_{\text{greedy}}$ در رأس $u$ یا $v$ اشتراک داشته باشد (در غیر این صورت الگوریتم حریصانه می‌توانست یال $(u, v)$ را به تطابق خود بیفزاید که با حریصانه بودن آن در تناقض است).
چون هر یال در $M_{\text{greedy}}$ حداکثر می‌تواند دو سر دو یال مختلف در $M_{\text{opt}}$ را لمس کند:
\begin{equation}
|M_{\text{opt}}| \le 2 \times |M_{\text{greedy}}| \implies \frac{|M_{\text{greedy}}|}{|M_{\text{opt}}|} \ge \frac{1}{2}
\end{equation}
بنابراین الگوریتم حریصانه همیشه حداقل نیمی از سود حالت ایده‌آل را کسب می‌کند.

\section{الگوریتم BALANCE با پیشنهادهای برابر}
در جهان واقعی، هر تبلیغ‌دهنده دارای یک \textbf{بودجه روزانه}\footnote{\lr{Daily Budget}} مشخص (مثلاً $B$ دلار) است. اگر پیشنهاد قیمت همه تبلیغ‌دهندگان برای کلمات کلیدی یکسان باشد (مثلاً ۱ دلار به ازای هر کلیک)، استراتژی حریصانه دلخواه ممکن است بودجه برخی تبلیغ‌دهندگان را زودتر از موعد تمام کرده و پرس‌وجوهای بعدی را بدون مشتری باقی بگذارد.

\begin{algorithmbox}[الگوریتم تعادل \enparen{BALANCE}]
هنگامی که پرس‌وجوی $q$ وارد می‌شود:
\begin{enumerate}
    \item تمامی تبلیغ‌دهندگانی را که برای $q$ پیشنهاد داده‌اند و هنوز بودجه آن‌ها تمام نشده است پیدا کنید.
    \item آگهی را به تبلیغ‌دهنده‌ای اختصاص دهید که \textbf{بیشترین بودجه باقیمانده} را دارد (در صورت تساوی، یکی را تصادفی انتخاب کنید).
\end{enumerate}
\end{algorithmbox}

\begin{theorem}[نسبت رقابتی الگوریتم BALANCE]
برای مسئله تخصیص آگهی با پیشنهادهای برابر و بودجه‌های متناهی، نسبت رقابتی الگوریتم BALANCE برابر است با:
\begin{equation}
\text{CR} = 1 - \frac{1}{e} \approx 0.632 \quad (63.2\%)
\end{equation}
این عدد نشان‌دهنده بهبود چشمگیر نسبت به الگوریتم حریصانه ساده $(50\%)$ است.
\end{theorem}

\section{مسئله عمومی \texorpdfstring{\lr{AdWords}}{AdWords} در گوگل (پیشنهادها و بودجه‌های نابرابر)}
در سازوکار واقعی سیستم Google AdWords:
\begin{itemize}
    \item هر تبلیغ‌دهنده $i$ برای پرس‌وجوی $q$ پیشنهاد قیمت متفاوتی به نام $b_{iq}$ ارائه می‌دهد.
    \item هر تبلیغ‌دهنده بودجه کل روزانه متفاوتی مانند $B_i$ دارد.
    \item هدف، بیشینه‌سازی کل درآمد ارزی کسب‌شده از کلیک‌هاست.
\end{itemize}

\subsection{تابع امتیازدهی عمومی AdWords}
اگر صرفاً به بیشترین پیشنهاد $(b_{iq})$ توجه کنیم، بودجه تبلیغ‌دهندگان بزرگ خیلی سریع مصرف شده و در ساعات پایانی روز درآمد افت می‌کند. از این رو، الگوریتم عمومی از یک تابع تنزیل غیرخطی استفاده می‌نماید.

فرض کنید $m_i$ مقدار پولی باشد که تبلیغ‌دهنده $i$ تاکنون خرج کرده است. کسر بودجه مصرف‌شده برابر با $f_i = m_i / B_i$ است.
هنگام ورود پرس‌وجوی $q$، تبلیغ‌دهنده برنده بر اساس \textbf{بیشینه امتیاز ترکیب‌شده} انتخاب می‌شود:
\begin{equation}
\Psi_i(q) = b_{iq} \times \left( 1 - e^{-(1 - f_i)} \right)
\end{equation}
\begin{itemize}
    \item اگر تبلیغ‌دهنده‌ای تقریباً تمام بودجه‌اش دست‌نخورده باشد $(f_i \approx 0)$، ضریب تصحیح نزدیک به $1 - 1/e \approx 0.63$ است.
    \item هرچه بودجه به اتمام نزدیک‌تر شود $(f_i \to 1)$، ضریب $1 - e^0 = 0$ شده و اولویت آن به شدت کاهش می‌یابد تا فضا برای سایرین باز شود.
\end{itemize}
اثبات شده است که این الگوریتم تعمیم‌یافته نیز نسبت رقابتی $1 - 1/e$ را تضمین می‌نماید.

\subsection{نرخ کلیک \texorpdfstring{\enparen{CTR}}{(CTR)} و رتبه‌بندی آگهی‌ها}
موتور جستجو تنها در صورتی درآمد کسب می‌کند که کاربر واقعاً بر روی آگهی کلیک نماید (مدل \lr{CPC}\footnote{\lr{Cost Per Click}}). بنابراین، ارزش انتظاری یک آگهی برابر است با:
\begin{equation}
\text{Expected Revenue} = \text{Bid} \times \text{CTR}
\end{equation}
که در آن \lr{CTR}\footnote{\lr{Click-Through Rate}} نسبت تعداد کلیک‌ها به تعداد نمایش آگهی است. آگهی‌ها بر اساس حاصل‌ضرب پیشنهاد در نرخ کلیک تاریخی (و نمره کیفیت متن آگهی) رتبه‌بندی می‌شوند.

\subsection{سازوکار حراج قیمت دوم \texorpdfstring{\enparen{Vickrey-Clarke-Groves Auction}}{(Vickrey-Clarke-Groves Auction)}}
گوگل برای ترغیب تبلیغ‌دهندگان به اعلام ارزش واقعی خود، از \textbf{حراج قیمت دوم تعمیم‌یافته}\footnote{\lr{Generalized Second-Price (GSP) Auction}} استفاده می‌کند:
برنده حراج، بالاترین پیشنهاد را ثبت کرده است، اما مبلغی که در ازای هر کلیک می‌پردازد برابر با پیشنهاد خودش نیست؛ بلکه برابر است با \textbf{حداقل مبلغی که برای شکست دادن پیشنهاد رقیب رتبه پایین‌تر لازم بوده است}.

\begin{theorem}[راست‌گویی به عنوان استراتژی مسلط در حراج قیمت دوم]
در حراج قیمت دوم، اعلام ارزش واقعی کالا توسط پیشنهاددهنده $(b = v)$ یک \textbf{استراتژی اکیداً مسلط}\footnote{\lr{Dominant Strategy}} است و هیچ عاملی نمی‌تواند با بیش‌برآورد یا کم‌برآورد سود خود را افزایش دهد.
\end{theorem}
\textbf{اثبات نظریه بازی‌ها:}
فرض کنید بالاترین پیشنهاد سایر رقبا برابر با $p$ باشد:
\begin{enumerate}
    \item اگر بازیکن بیش‌اظهاری کند $(b > v)$:
    تنها زمانی نتیجه تغییر می‌کند که پیشنهاد رقیب در بازه $v < p < b$ قرار گیرد. در این حالت، بازیکن برنده حراج می‌شود اما باید مبلغ $p$ را بپردازد که از ارزش واقعی کالا $(v)$ بیشتر است؛ یعنی سود خالص منفی $(v - p < 0)$ نصیبش می‌گردد.
    \item اگر بازیکن کم‌اظهاری کند $(b < v)$:
    تنها زمانی نتیجه تغییر می‌کند که پیشنهاد رقیب در بازه $b < p < v$ قرار گیرد. در این حالت، بازیکن فرصت برنده شدن در حراجی با سود خالص مثبت $(v - p > 0)$ را از دست می‌دهد.
\end{enumerate}
بنابراین، در تمامی حالات، ارائه پیشنهاد برابر با ارزش واقعی $(b = v)$ بهترین نتیجه ممکن را برای شرکت‌کننده تضمین می‌نماید.

\begin{examnote}
نکات پرکاربرد امتحانی در مبحث تبلیغات وب:
\begin{enumerate}
    \item کران پایینی الگوریتم حریصانه همیشه $1/2$ است و بدترین حالت آن زمانی رخ می‌دهد که یک تبلیغ‌دهنده مشترک میان دو پرس‌وجو، با اولین پرس‌وجو اشباع شود و پرس‌وجوی دوم بی‌مشتری بماند.
    \item عدد طلایی $1 - 1/e$ حد بالای تئوریک الگوریتم‌های برخط بدون اطلاع از توزیع ورودی‌ها است.
\end{enumerate}
\end{examnote}

\begin{summarybox}[جمع‌بندی فصل هشتم]
\begin{itemize}
    \item الگوریتم‌های برخط باید بدون آگاهی از رخدادهای آینده و به صورت آنی تصمیم‌گیری کنند.
    \item نسبت رقابتی عملکرد الگوریتم برخط را با سناریوی بهینه آفلاین مقایسه می‌نماید.
    \item الگوریتم حریصانه نسبت رقابتی $50\%$ و الگوریتم BALANCE نسبت رقابتی $63.2\%$ $(1-1/e)$ را تضمین می‌کنند.
    \item سیستم AdWords با ترکیب پیشنهادها، کسری بودجه، نرخ کلیک و حراج قیمت دوم، پایداری اقتصادی موتورهای جستجو را رقم می‌زند.
\end{itemize}
\end{summarybox}
